\chapter{Revisão Bibliográfica}
\label{sec:revisao}

\section{Considerações Iniciais}

A reconstrução de temperaturas a partir de redes de sensores envolve
relações entre as medições, disponibilidade dos registros e escolha das
entradas. Este capítulo apresenta os conceitos de monitoramento IoT,
séries temporais e sensoriamento virtual, seguidos dos métodos de
estimativa e das pesquisas que fundamentam a avaliação experimental.

\section{Conceitos e Técnicas Relevantes}
\label{sec:conceitos}

\subsection{Monitoramento Ambiental e Internet das Coisas}
\label{subsec:iot}

A Internet das Coisas (\textit{Internet of Things} -- IoT) integra
dispositivos de aquisição, comunicação e serviços que disponibilizam
medições para acompanhamento e análise. Em aplicações ambientais,
sensores distribuídos registram grandezas como temperatura e umidade,
produzindo séries associadas a diferentes posições físicas.

\citeonline{marques2026testbed} descrevem uma infraestrutura de laboratório
com dispositivos ESP32, comunicação multiprotocolo e armazenamento de
séries temporais. A disponibilidade de históricos dessa infraestrutura
permite investigar o aproveitamento conjunto das medições, além do
monitoramento direto de cada dispositivo. Contudo, a presença de um
registro no banco não assegura continuidade, sincronização ou
representatividade de todas as posições do ambiente.

\subsection{Séries Temporais e Pré-processamento de Dados}
\label{subsec:series_temporais}

Uma série temporal reúne observações associadas a instantes ou
intervalos de tempo. Sensores podem registrar valores em instantes
diferentes e apresentar lacunas, mesmo quando possuem frequências
nominais semelhantes. A análise conjunta exige verificar identificadores,
convenções temporais, duplicatas, valores inválidos e disponibilidade
simultânea.

Uma estratégia de sincronização é agregar as medições em intervalos
regulares. Para o sensor $s$ e a janela $t$, a média é
\begin{equation}
    \overline{T}_{s,t}=\frac{1}{n_{s,t}}
    \sum_{j=1}^{n_{s,t}}T_{s,t,j},
    \label{eq:media_temporal}
\end{equation}
em que $n_{s,t}>0$ é a quantidade de medições disponíveis e $T_{s,t,j}$
a $j$-ésima temperatura da janela. Intervalos maiores reduzem a resolução
temporal e podem facilitar a coincidência dos registros, enquanto
intervalos menores preservam variações mais rápidas.

Valores ausentes podem ser preservados, excluídos ou imputados, conforme
o objetivo e as características dos dados. \citeonline{alsalehy2025timeseries}
discutem a importância de considerar a estrutura temporal na detecção
de anomalias e no preenchimento de lacunas ambientais. Qualquer regra
adotada deve ser documentada para permitir interpretar suas consequências
sobre o conjunto efetivamente avaliado.

\subsection{Aprendizado de Máquina e Enriquecimento dos Dados}
\label{subsec:ia_enriquecimento}

O aprendizado de máquina é a abordagem de inteligência artificial (IA)
empregada neste estudo. No aprendizado supervisionado, observações
históricas relacionam variáveis de entrada a uma resposta conhecida,
permitindo ajustar um modelo e avaliar suas estimativas em outras
observações. A Scikit-learn disponibiliza implementações para esse
procedimento \cite{pedregosa2011scikit}; entre os métodos estudados,
a Random Forest combina árvores de decisão \cite{breiman2001random}.

O enriquecimento amplia as informações associadas aos registros.
No contexto deste trabalho, compreende vincular as medições a metadados
dos dispositivos e às posições físicas, além de combinar temperatura
e umidade em uma mesma janela temporal. As informações espaciais podem
orientar a seleção de entradas, enquanto a inclusão de outra grandeza
ambiental permite investigar sua contribuição preditiva. O benefício
precisa ser demonstrado na avaliação, pois aumentar as informações
não assegura menor erro.

A integração de grandezas ambientais e registros de energia também
é explorada por \citeonline{tavoni2026digitaltwin}, que utilizam essas
informações na previsão de consumo em uma arquitetura de gêmeo digital.
Essa relação oferece contexto para o enriquecimento investigado aqui,
embora a variável de resposta e o sistema desenvolvido sejam distintos.

\subsection{Sensoriamento Virtual de Temperatura}
\label{subsec:sensoriamento_virtual}

Um sensor virtual utiliza informações disponíveis para estimar uma
grandeza que poderia ser medida por um sensor físico. Na reconstrução
de temperaturas, o problema pode ser representado por
\begin{equation}
    \widehat{T}_p(t)=f\bigl(T_1(t),\ldots,T_k(t)\bigr),
    \label{eq:sensoriamento_virtual}
\end{equation}
em que $p$ é o alvo, $k$ é a quantidade de entradas e $f$ é o método
de estimativa. As medições do próprio alvo não podem integrar essas
entradas durante a reconstrução; seu histórico pode fornecer a resposta
para treinamento supervisionado.

\citeonline{brunello2021virtual} investigam essa abordagem em um escritório
com sensores distribuídos irregularmente, combinando seleção de entradas
e avaliação de preditores. A distância física e a correlação entre
séries são critérios distintos: o primeiro utiliza coordenadas, enquanto
o segundo expressa associação entre observações. Mais entradas podem
adicionar informação ou redundância, tornando necessário avaliar seu
efeito em dados separados do ajuste e da seleção.

\subsection{Métodos de Estimativa e Avaliação}
\label{subsec:metodos_estimativa}

A média aritmética constitui uma referência sem treinamento:
\begin{equation}
    \widehat{T}_p(t)=\frac{1}{k}\sum_{i=1}^{k}T_i(t).
    \label{eq:media_sensores}
\end{equation}
Ela atribui pesos iguais às entradas e não ajusta diferenças sistemáticas
entre o nível térmico dos sensores e o do alvo.

A interpolação pelo inverso da distância (\textit{Inverse Distance
Weighting} -- IDW) atribui maior peso a observações próximas:
\begin{equation}
    \widehat{T}_p(t)=
    \frac{\sum_{i=1}^{k}T_i(t)d_{ip}^{-q}}
         {\sum_{i=1}^{k}d_{ip}^{-q}},
    \label{eq:idw}
\end{equation}
em que $d_{ip}>0$ é a distância ao ponto estimado e $q>0$ controla
o decaimento dos pesos. Na coincidência espacial com um sensor,
atribui-se seu valor diretamente. \citeonline{jedermann2011interpolation}
comparam IDW e Kriging na reconstrução de perfis de temperatura.
A interpolação fornece valores entre pontos observados, mas sua
representação gráfica não substitui a validação contra medições.

A regressão linear estima o alvo por uma combinação ajustada das entradas:
\begin{equation}
    \widehat{T}_p(t)=\beta_0+\sum_{i=1}^{k}\beta_iT_i(t),
    \label{eq:regressao_linear}
\end{equation}
em que $\beta_0$ é o intercepto e $\beta_i$ são os coeficientes estimados
no treinamento. A Random Forest, apresentada por
\citeonline{breiman2001random}, combina árvores de decisão; em regressão,
as estimativas das árvores são agregadas. Essa estrutura permite
representar relações não lineares, mas não assegura menor erro que
uma alternativa linear em qualquer conjunto de dados.

A comparação utiliza o erro absoluto médio (\textit{Mean Absolute
Error} -- MAE) e a raiz do erro quadrático médio (\textit{Root Mean
Squared Error} -- RMSE):
\begin{equation}
    \mathrm{MAE}=\frac{1}{n}\sum_{j=1}^{n}|T_j-\widehat{T}_j|,
    \label{eq:mae}
\end{equation}
\begin{equation}
    \mathrm{RMSE}=\sqrt{\frac{1}{n}
    \sum_{j=1}^{n}(T_j-\widehat{T}_j)^2}.
    \label{eq:rmse}
\end{equation}
As duas métricas possuem a unidade da temperatura, mas o RMSE penaliza
mais os erros elevados. $T_j$ é o valor observado, $\widehat{T}_j$ é a
estimativa e $n$ é o número de observações. Sua interpretação exige
informar a divisão temporal, os alvos e as amostras utilizados.

\subsection{Relação com o Conforto Térmico}
\label{subsec:conforto_termico}

Temperatura e umidade são relevantes à caracterização ambiental,
mas a interpretação do conforto envolve outras condições. A apresentação
institucional da ASHRAE Standard 55 relaciona fatores ambientais
--- temperatura, radiação térmica, umidade e velocidade do ar ---
e pessoais, como atividade e vestimenta \cite{ashrae55overview}.
Assim, uma temperatura reconstruída com baixo erro não permite,
isoladamente, classificar a satisfação térmica dos ocupantes.

O presente conjunto não inclui caracterização da radiação térmica,
da velocidade do ar, da atividade ou da vestimenta, nem votos de
sensação térmica. Esses fatores não foram substituídos por valores
arbitrários para produzir um índice. A reconstrução e a representação
espacial são tratadas como recursos analíticos que poderiam apoiar
uma avaliação de conforto futura, mediante protocolo e dados adequados.

\section{Trabalhos Relacionados}
\label{sec:trabalhos_relacionados}

\citeonline{marques2026testbed} apresentam uma infraestrutura IoT com
ESP32, MQTT, CoAP, atualização remota de firmware e observabilidade.
Os autores relatam redução de 55,7\% no consumo mensal de energia em
seu cenário de controle de climatização, mantendo condições de conforto.
Essa evidência caracteriza a aplicação da infraestrutura estudada
pelos autores; o presente trabalho utiliza históricos de temperatura
e não reproduz o experimento energético.

\citeonline{tavoni2026digitaltwin} apresentam uma arquitetura modular
de gêmeo digital para otimização energética em um laboratório da USP.
A solução integra sensores de temperatura e umidade DHT22, sensores
de corrente SCT-013, comunicação MQTT, Eclipse Ditto e uma interface
em Unity. O componente preditivo utiliza XGBoost para estimar consumo
em kWh, com valores anteriores de energia, temperatura e umidade,
validação por \textit{TimeSeriesSplit} e parada antecipada no treinamento.

O artigo descreve cenários de simulação com diferentes configurações
de operação e reconhece limitações relacionadas ao período experimental
e à generalização. Sua contribuição para este TCC está na integração
de dados heterogêneos e no uso de aprendizado de máquina para apoiar
análises da climatização. Neste trabalho, a resposta é a temperatura
de um sensor e a avaliação usa uma divisão temporal fixa; não foram
implementados gêmeo digital, XGBoost, simulação de consumo ou atuação.
O SCT-013 descrito pelos autores também difere do modelo ACS712
identificado na extração aqui analisada. Portanto, a unidade kWh
reportada naquele estudo não confirma a unidade do campo
\texttt{energy} deste conjunto.

\citeonline{brunello2021virtual} constituem a referência mais próxima
do problema investigado. Os autores analisam métricas de distância
espaciais e estatísticas, desenvolvem uma combinação por programação
genética e selecionam subconjuntos de sensores para estimar temperaturas
em um escritório. Avaliam diferentes estratégias de predição e a
confiabilidade das saídas. O presente estudo adota uma comparação mais
restrita, com métodos simples e duas regras de seleção, aplicada aos
dispositivos e históricos disponíveis no LASDPC.

\citeonline{yuan2025optimization} investigam reconstrução de campos
térmicos em regimes estacionários e transientes, com cenários numéricos
de condução de calor. Comparam configurações esparsas de sensores e
métodos baseados em decomposição ortogonal própria (\textit{Proper
Orthogonal Decomposition} -- POD), incluindo Gappy POD e Gappy C-POD.
O estudo fundamenta a relevância da configuração das medições. Entretanto,
a otimização de posições em cenários numéricos difere da condição deste
TCC, que utiliza sensores reais previamente instalados.

\citeonline{jedermann2011interpolation} aplicam interpolação espacial
a dados de armazenamento e transporte de alimentos, comparando Kriging
com ponderação pelo inverso da distância. Nos conjuntos avaliados pelos
autores, o erro de Kriging foi aproximadamente 20\% menor que o da
referência IDW. O trabalho sustenta a discussão de relações espaciais,
mas não transfere esse resultado diretamente ao ambiente do LASDPC.
Neste TCC, a IDW é usada apenas para ilustrar as médias espaciais.

\citeonline{alsalehy2025timeseries} avaliam identificação de anomalias e
imputação em séries de gases e variáveis meteorológicas, combinando
métodos estatísticos e algoritmos de detecção, além de interpolação
linear, PCHIP e Akima. Sua relação com o presente estudo é a atenção
à continuidade e à qualidade das séries. A decisão experimental aqui
adotada foi selecionar janelas completas, sem reproduzir as técnicas
de imputação desse artigo.

Por fim, \citeonline{cho2022preprocessing} investigam combinações de
tratamento de valores ausentes e desbalanceamento para classificação
industrial. Embora a aplicação seja distinta da reconstrução de
temperaturas, o artigo reforça que a preparação dos dados integra
o procedimento de avaliação e deve ser descrita com precisão.

\section{Considerações Finais}
\label{sec:consideracoes_finais_revisao}

Os trabalhos sobre infraestrutura e gêmeos digitais contextualizam
a integração das medições e suas aplicações na climatização.
A literatura relaciona o sensoriamento virtual à seleção de entradas,
aos métodos de estimativa e à qualidade das medições. Este trabalho
reúne essas perspectivas em uma avaliação delimitada de dados reais,
sem propor um novo algoritmo ou otimizar a posição física dos sensores.
O capítulo seguinte descreve a infraestrutura, os dados e o protocolo,
e apresenta os resultados nas condições efetivamente disponíveis.
